\section{Related Work}
\label{sec:related}

\emph{Lie detection.} Models lie when instructed to, and instruction pairs (the same statement produced
under an honest and under a deceptive instruction) are the standard training data for deception probes
\cite{zou2023representation,goldowskydill2025detecting}. Black-box detectors use only inputs and
outputs: Pacchiardi et al.~\cite{pacchiardi2023catch} ask unrelated follow-up questions after a
suspected lie and classify the answers, and Liars' Bench adds an LLM judge that reads the transcript
\cite{kretschmar2025liars}. The judge is the strongest overall but fails on \hpkr, one of the settings
where ``it's not possible to determine whether the model lied from the transcript alone'', and each
black-box check costs at least one additional model call. Goldowsky-Dill et
al.~\cite{goldowskydill2025detecting} trained logistic-regression probes on Llama-3.3-70B-Instruct
and reached AUROCs of 0.96--0.999 on the model's own deception in insider trading and sandbagging, yet
their probe also responds to whether ``the scenario [is] related to deception'' and to morality.

\emph{Why probes fail to transfer.} A probe that is nearly perfect in distribution can be near or below
chance elsewhere. With fewer examples than the dimensions of the residual stream, many directions fit
a training set equally well \cite{damour2022underspecification}, and logistic regression chooses among
them by the margin of the particular training set \cite{soudry2018implicit}. Discriminative learning
uses any feature that suffices on a given dataset \cite{geirhos2020shortcut}: probes have latched onto
negation and polarity \cite{levinstein2025still,buerger2024truth}, text probability
\cite{marks2023geometry}, deception-related scenarios \cite{goldowskydill2025detecting} and style
\cite{kumar2026pressure}. Off-policy training data harm ``text-ambiguous'' behaviours such as deception
\cite{kirch2025offpolicy}, and deception is heterogeneous: probes matched to a type of deception
outperform a universal one \cite{natarajan2026targeted}. Probes fall \emph{below} chance where a
correlate of the label reverses, as on negated facts \cite{levinstein2025still}, on
expectation-inverted lying \cite{ying2026truthfulness} and on \hpkr{}
\cite{kretschmar2025liars,moustafa2026beyond}. Remedies constrain the direction: to the difference of
means \cite{marks2023geometry}, to selected principal components \cite{yoo2026subspace}, or to persona
coordinates \cite{mahadik2026persona}. We take the last route.

\emph{Liars' Bench.} The benchmark \cite{kretschmar2025liars} collects \V{lb-examples} on-policy lies
and honest responses across seven datasets and four models, classified by the object of belief and the
reason for lying. It evaluates a \emph{Mean Probe} trained on RepE instruction pairs
\cite{zou2023representation}, a \emph{Follow-Up Probe} read at a prefilled ``No'' after the question
``Was your last response deceptive?'', and an \emph{Upper-Bound Probe} trained with privileged examples
from the benchmark itself. On \hpkr, built from the hazardous-knowledge questions of WMDP
\cite{li2024wmdp}, the judge's recall is almost zero and the Mean and Follow-Up Probes invert
(\V{lb-mean-auroc-hpkr} and \V{lb-followup-auroc-hpkr} AUROC). Only the Upper-Bound Probe, trained
with labelled \hpkr{} examples, detects them reliably.

\emph{Personas as directions.} Persona vectors \cite{chen2025persona} are activation directions for
character traits, used to monitor and steer them. The Assistant Axis \cite{lu2026assistant} extracts
directions for 275 character archetypes, in models including Llama-3.3-70B. Mahadik and
Skapars~\cite{mahadik2026persona} turned this machinery into a monitor for deception and sycophancy.
They build a deception axis from eight deceptive and seven honest personas and a sycophancy axis from
five sycophantic and five anti-sycophantic ones, and use the contrast directions zero-shot and the top
three principal components as features for probes. On Llama~3.2-3B and Llama~3-8B these persona
coordinates transfer better than probes on raw activations, across datasets that do not include
\hpkr. The same literature marks the limits of the approach: roles differ in style as well as in
honesty \cite{joshi2024personas}, and a model that switches to an ignorant persona deceives although
``no persona is ever lying'' \cite{marks2026persona}, which is the shape of a denial of knowledge.

\emph{Reading what a model knows.} Several lines of work show that a model's internal states record
whether it knows an answer. Language models can be trained to predict whether they can answer a
question correctly \cite{kadavath2022language}; a probe separates answerable from unanswerable
questions even when the model goes on to invent an answer \cite{slobodkin2023curious}; how much a model
knows about an entity can be predicted from its representation alone \cite{gottesman2024estimating};
and sparse-autoencoder directions detect whether the model recognises an entity, and steer its
refusals \cite{ferrando2025know}. Our use differs in three ways. The direction is read at the tokens of
a denial written for the model, so the words are fixed and only the knowledge varies. It is a
difference of means over questions marked as known or unknowable, not a classifier trained on the
model's correctness. And it is used to detect lies, not to predict hallucination.
